
This work builds on SMT-guided program repair, neural code generation, and static analysis for typed languages. The hypothesis---incremental SMT for LLM code repair in typed Python---lies at their intersection, though empirical validation was blocked.

\subsubsection{SMT-Guided Program Repair}

Angelix and Prophet apply SMT solvers to synthesize patches for buggy programs by encoding correctness specifications as constraints. These systems work on small human-written patches (1--10 lines) and assume localized modifications. Constraint extraction relies on predefined templates for common bug patterns. While effective for human code repair, they do not scale to full LLM-generated programs where bug locations are unknown and programs exceed 100 lines.

The hypothesis proposed in this work extends this paradigm to LLM code repair workflows: instead of batch re-verification on each iteration, incremental SMT re-verifies only modified functions and dependencies. The key differences are scale (full programs, not patches) and code source (LLM-generated code may have different quality characteristics). However, transferability remains untested due to validation failure.

\subsubsection{Neural Code Generation}

AlphaCode and CodeT5 demonstrate that large-scale pretraining enables LLMs to generate functionally correct code on benchmarks. These systems rely on test-suite validation: generate multiple candidates, filter by passing tests, rank by confidence. While effective for programming competitions, test suites provide incomplete correctness guarantees.

Formal verification could address this gap by providing SMT-based correctness proofs. However, existing neural code generation systems do not integrate static analysis or SMT solvers. This work proposed adding incremental SMT verification to LLM repair loops, but could not validate whether LLM-generated code has extractable constraints or whether SMT provides practical benefit over test suites.

\subsubsection{Static Analysis for Typed Languages}

Prusti (Rust) and Pyre (Python) extract SMT constraints from type annotations and function contracts. These tools demonstrate that type annotations map to first-order logic predicates suitable for SMT solving. Prusti uses Rust ownership types and developer-written specifications; Pyre uses Python type hints and optional runtime assertions.

Prior work validates these analyzers on human-written typed code. Whether LLM-generated code has sufficient annotation quality for constraint extraction is an open question. This work hypothesized $\geq 90$\% extraction success, but could not test this due to API authentication failures. The constraint extraction pipeline was validated on error files, confirming architectural soundness on negative cases, but the core claim remains empirically unverified.

\subsubsection{Neural Program Repair}

CURE and CoCoNut demonstrate that neural program repair typically modifies 1--3 lines in human code (80\% of repairs are localized). The hypothesis assumed LLM repairs exhibit similar locality. If true, dependency cones remain small (<30\%), enabling incremental SMT speedup. However, LLM repair patterns may differ from human edits---broader modifications, cross-module dependencies---which could invalidate speedup claims. Repair locality could not be tested due to prerequisite failures.

\subsubsection{Benchmarks for Typed Code}

HumanEval provides 164 Python programming problems with test suites but lacks type annotations required for static analysis. APPS and CodeContests similarly focus on functional correctness without formal specifications. This work fills this gap by mechanically generating 100 type-annotated prompts. This artifact is reusable for future constraint extraction research, independent of the validation failure.

\subsubsection{Positioning}

This hypothesis occupies unexplored territory: combining SMT-guided repair (established for human patches) with neural code generation (lacking verification) in typed Python with Pydantic. The key novelty is incremental SMT for full LLM programs, but this remains a proposal pending validation. The verified contribution is the dataset artifact and pipeline architecture.
